% Abstract (all). Full-width paragraph on body page 1. Keep under ~230 words.
\textbf{Abstract.} Every AI model runs in a data centre that draws electricity, evaporates water for cooling and is
built from mined hardware that ends as e-waste. Data centres used about 1.5\% of the world's electricity in 2024, and
the IEA expects that to roughly double by 2030. We place this growth in three phases since 1950 (the Great
Acceleration, the cloud era and the AI era) and ask why cloud growth left data-centre electricity almost flat while AI
growth does not. Combining Epoch AI's 96 AI data centres (15.1\,GW today, 50.3\,GW planned by 2030) with WRI Aqueduct,
the GHSL population grid, Ember, USGS and published energy series, we find that AI training compute broke upward in
2011, but US data-centre electricity only accelerated after 2018 (from about 7\% to 18\% a year), once building
efficiency had stopped improving. Weighted by power, 34\% of tracked AI capacity sits in basins of high or extremely
high water stress today and 39--45\% of the 2030 plan; 21 million people live within 25\,km of those sites. The sites
consume an estimated 143--624 billion litres of water a year, 67--86\% of it at the power plants that supply them.
Against the global freshwater boundary this is under 0.1\%; in Google's Mesa basin, withdrawals already exceed supply in
all twelve months. At the planned Visakhapatnam hub, water risk shaped the cooling design but was not a siting test, and no binding framework required one. Capacity is far more concentrated than people (country-level Gini 0.94), and the minerals come from a
few producing countries. AI's footprint is small globally and concentrated locally, so responses have to be matched to
the local, regional and global scale of each impact.
